% from ega4-2-fr.tex
\oldpage[IV-2]{64}

géométrique au-dessus d'un corps algébriquement clos $\Omega$~; la donnée d'un
«~corps de définition~» $k$ plus petit (i.e., d'un schéma sur $k$ définissant
$V$ par extension du corps de base à $\Omega$) est pour lui une structure
supplémentaire et, dans une certaine mesure, secondaire, de sorte que sa
qualification d'«~absolu~» ou «~intrinsèque~» d'un côté, «~relatif à $k$~» de
l'autre, est à l'opposé de la nôtre\footnote{Le point de vue de Zariski est
déjà plus proche du nôtre, et sa terminologie n'est généralement pas en conflit
avec celle introduite ici.}. Nous éviterons donc par la suite d'utiliser le
qualificatif «~absolu~», qui peut prêter à confusion, et lorsque nous
utiliserons la terminologie abrégée introduite plus haut (en conflit avec la
terminologie reçue), nous renverrons à (4.5.12) pour éviter toute ambiguïté.

\item[(ii)] La proposition (4.5.9) donne un critère «~birationnel~» (autrement
dit, ne dépendant que du corps résiduel au point générique) pour qu'un
$k$-préschéma $X$ soit géométriquement irréductible. Il n'y a pas de critère
analogue pour que $X$ soit géométriquement connexe. Par exemple, prenons
$X=\operatorname{Spec}(\mathbf R[[T,U]]/(T^2+U^2))$ ($T,U$ indéterminées)~;
$X$ est un $\mathbf R$-schéma intègre, et son corps de fonctions rationnelles
$k(x)$ est $\mathbf C((T))$, comme on le vérifie aisément~; $X\otimes_{\mathbf
R}\mathbf C$ se décompose en deux composantes irréductibles (les deux
«~droites isotropes~») qui ont un point commun (l'idéal maximal image dans
$\mathbf C[[T,U]]/(T^2+U^2)$ de l'idéal maximal $(T)+(U)$ de
$\mathbf C[[T,U]]$)~; donc $X\otimes_{\mathbf R}\mathbf C$ est connexe. Soit
$X'$ le complémentaire (ouvert) du point fermé de $X$ correspondant à l'idéal
maximal $(T)+(U)$ de $\mathbf R[[T,U]]$~; alors $X'\otimes_{\mathbf R}\mathbf C$
n'est pas connexe, bien que $X$ et $X'$ aient même corps de fonctions
rationnelles.
\end{enumerate}
\end{remark}

\begin{proposition}[4.5.13]
\label{IV.2.4.5.13-fr}
Soient $X$, $Y$ deux $k$-préschémas, $f:Y\to X$ un $k$-morphisme. Supposons
$Y$ géométriquement connexe et non vide, et $X$ connexe. Alors $X$ est
géométriquement connexe.
\end{proposition}

\begin{proof}
Soient $\Omega$ une clôture algébrique de $k$, $X'=X_{(\Omega)}$,
$Y'=Y_{(\Omega)}$, $f'=f_{(\Omega)}:Y'\to X'$~; il faut montrer que $X'$ est
connexe. Soit $U'$ une partie ouverte et fermée non vide de $X'$~; notons que
les morphismes $p:X'\to X$ et $q:Y'\to Y$ sont \emph{ouverts} (2.4.10) et
\emph{fermés} puisque $\Omega$ est extension algébrique de $k$ (II, 6.1.10).
Donc $U=p(U')$ est ouvert et fermé non vide dans $X$, et par suite égal à $X$.
Comme $Y$ est non vide, on en conclut (I, 3.4.7) que
$V'=f'^{-1}(U')$ n'est pas vide~; d'ailleurs, c'est une partie ouverte et
fermée de $Y'$, et ce dernier espace est connexe par hypothèse~; donc $V'=Y'$.
On en déduit que $U'=X'$, sans quoi le même raisonnement appliqué à l'ensemble
ouvert et fermé $X'-U'$, montrerait que $f'^{-1}(X'-U')=Y'$, ce qui est absurde
puisque $Y'$ n'est pas vide.
\end{proof}

\begin{corollary}[4.5.13.1]
\label{IV.2.4.5.13.1-fr}
\begin{enumerate}
\item[(i)] Soient $X$, $Y$ deux $k$-préschémas, $f:Y\to X$ un $k$-morphisme.
Si $Y$ est géométriquement connexe et non vide, la composante connexe $X_0$ de
$X$ contenant $f(Y)$ est géométriquement connexe.
\item[(ii)] Soit $X$ un $k$-préschéma. Si $Y$ est une composante irréductible
de $X$ qui est géométriquement irréductible, alors la composante connexe $X_0$
de $X$ contenant $Y$ est géométriquement connexe.
\end{enumerate}
\end{corollary}

\begin{proof}
Pour (i), on peut supposer $Y$ réduit, de sorte que $f$ se factorise en
$Y\xrightarrow{g}X_0\xrightarrow{j}X$, $X_0$

